% Part: model-theory
% Chapter: interpolation
% Section: definability

\documentclass[../../../include/open-logic-section]{subfiles}

\begin{document}

\olfileid{mod}{int}{def}

\olsection{Teorema Definabilitas}

Salah siji panganggone teorema interpolasi kang wigati yaiku
teorema definabilitas Beth. Kanggo netepake relasi
$n$-papan~$R$, kita bisa menehi !!a{formula}~$!C$
kanthi $n$ !!{variable}s bebas kang ora ngemot~$R$.
Iki dadi definisi \emph{eksplisit} kanggo~$R$ adhedhasar~$!C$.
Kita uga bisa kandha yen sawijining teori~$\Sigma(P)$
ing !!{language} kang ngemot !!{predicate}
$n$-papan~$P$ netepake~$P$ kanthi eksplisit yen teori iku
ngemot (utawa paling ora nduweni akibat) definisi
eksplisit kang diformalkan, yaiku
\[
\Sigma(P) \Entails \lforall[x_1][\dots
  \lforall[x_n][(\Atom{P}{x_1,\dots, x_n} \liff !C(x_1, \dots,
    x_n))]].
\]
Nanging definisi eksplisit mung siji cara kanggo
netepake relasi kanthi tuntas. Sawijining teori
uga bisa nemtokake interpretasi~$P$ kanthi tuntas
lumantar interpretasi bagean !!{language} liyane
ing saben model. Teorema definabilitas nyatakake:
saben teori kang nemtokake interpretasi~$P$
kanthi cara iki---yaiku \emph{netepake~$P$ kanthi implisit}---
uga netepake kanthi eksplisit.

\begin{defn}
Upamane $\Lang{L}$ iku !!a{language} kang ora ngemot
!!{predicate}~$P$. Himpunan $\Sigma(P)$ saka
!!{sentence}s ing $\Lang{L} \cup \{P\}$
\emph{netepake~$P$ kanthi eksplisit} yen lan mung yen
ana !!a{formula}~$!C(x_1, \dots, x_n)$ ing $\Lang{L}$
saengga
\[
\Sigma(P) \Entails \lforall[x_1][\dots
  \lforall[x_n][(\Atom{P}{x_1,\dots, x_n} \liff !C(x_1, \dots,
    x_n))]].
\]
\end{defn}

\begin{defn}
Upamane $\Lang{L}$ iku !!a{language} kang ora ngemot
!!{predicate}s~$P$ lan~$P'$. Himpunan $\Sigma(P)$
saka !!{sentence}s ing $\Lang{L} \cup \{P\}$
\emph{netepake $P$ kanthi implisit} yen lan mung yen
\[
\Sigma(P) \cup \Sigma(P') \Entails \lforall[x_1][\dots
  \lforall[x_n][(\Atom{P}{x_1,\dots, x_n} \liff \Atom{P'}{x_1,\dots,
      x_n})]],
\]
ing ngendi $\Sigma(P')$ dipikolehi saka panggantos
seragam $P$ dadi $P'$ ing $\Sigma(P)$.
\end{defn}

Tegese, kanggo saben model $\Struct{M}$ lan $R, R' \subseteq
\Domain{M}^n$, yen $\Expan{M}{R} \Entails \Sigma(P)$ lan
$\Expan{M}{R'} \Entails \Sigma(P')$, mula $R=R'$.
Ing kene $\Expan{M}{R}$ iku !!{structure}~$\Struct{M'}$
kanggo ekspansi saka $\Lang{L}$ menyang
$\Lang{L} \cup \{P\}$ kang nduweni
$\Assign{P}{M'} = R$; semono uga kanggo
$\Expan{M}{R'}$ kanthi lambang predikat kang wis diganti.

\begin{thm}[Teorema Definabilitas Beth]
Himpunan $\Sigma(P)$ saka !!{sentence}s ing $\Lang{L}
  \cup\{P\}$ netepake $P$ kanthi implisit yen lan mung yen
  $\Sigma(P)$ netepake $P$ kanthi eksplisit.
Cathetan koreksi sumber OLPL-120:
teorema iki nganggo himpunan ukara kaya rong
definisi sadurunge, dudu himpunan formula bebas.
\end{thm}

\begin{proof}
Yen $\Sigma(P)$ netepake $P$ kanthi eksplisit, mula
loro pratelan iki bener:
\begin{align*}
  \Sigma(P) & \Entails & \lforall[x_1][\dots \lforall[x_n]
    [(\Atom{P}{x_1,\dots, x_n} \liff !C(x_1,\dots,x_n))]]\\
  \Sigma(P') & \Entails & \lforall[x_1][\dots \lforall[x_n]
    [(\Atom{P'}{x_1,\dots, x_n} \liff !C(x_1,\dots,x_n))]]
\end{align*}
lan kesimpulane banjur langsung tumut.
Kanggo arah kosok baline, upamane $\Sigma(P)$
netepake $P$ kanthi implisit. Dhisik kita tambah
!!{constant}s $c_1$, \dots,~$c_n$ menyang $\Lang{L}$.
Banjur
\[
\Sigma(P) \cup \Sigma(P') \Entails
\Atom{P}{c_1, \dots, c_n} \to  \Atom{P'}{c_1, \dots, c_n}.
\]
Miturut kompaktitas, ana himpunan winates
$\Delta_0 \subseteq \Sigma(P)$ lan
$\Delta_1 \subseteq \Sigma(P')$ saengga
\[
\Delta_0 \cup \Delta_1 \Entails
\Atom{P}{c_1, \dots, c_n} \to \Atom{P'}{c_1, \dots, c_n}.
\]
Sebut $!D(P)$ konjungsi kabeh !!{sentence}s $!A(P)$
kang $!A(P) \in \Delta_0$ utawa $!A(P') \in \Delta_1$;
sebut $!D(P')$ konjungsi kabeh !!{sentence}s $!A(P')$
kang $!A(P) \in \Delta_0$ utawa $!A(P') \in \Delta_1$.
Banjur $!D(P) \land !D(P') \Entails
\Atom{P}{c_1, \dots, c_n} \to \Atom{P'}{c_1,\dots,c_n}$.
Cathetan koreksi sumber OLPL-119:
predikat ing sisih tengen ditulis nganggo wangun atomik
kang padha karo premis lan rumus sabanjure.
Kita bisa nata maneh supaya saben !!{predicate}
katon ing sisih dhewe-dhewe saka $\Entails$:
\[
!D(P) \land \Atom{P}{c_1, \dots, c_n} \Entails
!D(P') \to \Atom{P'}{c_1, \dots, c_n}.
\]
Miturut Teorema Interpolasi Craig, ana !!a{sentence}
$!C(c_1,\dots, c_n)$ kang ora ngemot $P$ utawa $P'$
saengga:
\begin{align*}
  !D(P) \land \Atom{P}{c_1, \dots, c_n} & \Entails !C(c_1,\dots, c_n); \\
  !C(c_1,\dots, c_n) & \Entails !D(P') \to \Atom{P'}{c_1, \dots, c_n}.
\end{align*}
Saka akibat kang kapisan, kita oleh $!D(P) \Entails
\Atom{P}{c_1,\dots, c_n} \lif !C(c_1,\dots, c_n)$.
Saka akibat kang kapindho, amarga model
$\Lang{L} \cup \{P\}$ yaiku $\Expan{M}{R}
\Entails !A(P)$ yen lan mung yen model
$\Lang{L} \cup \{P'\}$ kang cocog,
$\Expan{M}{R} \models !A(P')$, kita oleh
$!C(c_1,\dots, c_n) \Entails !D(P) \lif
\Atom{P}{c_1,\dots, c_n}$, mula:
\[
!D(P) \Entails !C(c_1,\dots,c_n) \to \Atom{P}{c_1,\dots, c_n}.
\]
Yen loro asil iki digandhengake,
$!D(P) \Entails \Atom{P}{c_1,\dots, c_n}
\liff !C(c_1, \dots, c_n)$; miturut monotonisitas
lan generalisasi, kita uga oleh
\[
\Sigma(P) \Entails
\lforall[x_1][\dots\lforall[x_n][(\Atom{P}{x_1,\dots, x_n} \liff
    !C(x_1,\dots, x_n))]].
\]
\end{proof}

\end{document}
